\usepackage{enumitem}

% Очищуємо зайві поля
\faculty{}
\discipline{}
\course{}
\group{}
\speciality{}
\leader{}
\degree{}
\position{}
\reviewer{}
\controller{}
\approver{}

% --- ДАНІ КНИГИ ---
% Автор
\newcommand{\bookauthor}{Сергій Євдокимов}

% Назва (для використання в метаданих)
\newcommand{\booktitle}{Протидія шахрайству: Платформи оголошень та цифровi псевдопослуги}

% Місто
\newcommand{\bookcity}{Дніпро}

% Рік
\newcommand{\bookyear}{2026}

% --- ДОДАТКОВІ ПОЛЯ ДЛЯ ДРУГОЇ СТОРІНКИ ТИТУЛКИ ---
% Основні метадані для другої титульної сторінки — встановлені тут
\newcommand{\udc}{УДК 004.056.5:343.533(075.8)}
\newcommand{\authorslist}{Євдокимов С. Г.}
\newcommand{\recommendation}{Рекомендовано до друку Вченою радою [Назва установи] (протокол № [X] від [Дата] 2026 р.)}
\newcommand{\consultants}{[Прізвище І. П.], науковий ступінь, вчене звання.}
\newcommand{\reviewers}{%
	[Прізвище І. П.], науковий ступінь, вчене звання; \\
	[Прізвище І. П.], науковий ступінь, вчене звання.%
}

\newcommand{\biblioentry}{Євдокимов С. Г. Протидія шахрайству: Платформи оголошень та цифровi псевдопослуги. — Дніпро, 2026. — 202 с. : іл.}

\newcommand{\annotationshort}{%
	\textbf{Анотація.} У книзі системно викладено теоретичні та практичні засади протидії сучасним цифровим загрозам. 
	Автор аналізує схеми шахрайства на платформах оголошень, у соцмережах та месенджерах, приділяючи особливу увагу 
	соціальній інженерії та фішингу. Видання містить алгоритми перевірки контрагентів та технічні поради щодо 
	захисту персональних даних. Розраховано на широке коло користувачів, які прагнуть підвищити рівень 
	цифрової грамотності та безпеки.
}

\newcommand{\copyleft}{© Євдокимов С. Г., 2026 \\ © Дизайн та оформлення, 2026}
